\answer{ex:08:1}{(а)~$\approx1.7\cdot10^6$. (б)~$g=\frac{0.9}{\sqrt{0.19}}\,\sigma_{\text{після}}/\sigma_{\text{до}}\approx16.5$: відповідь залежить лише від відношення шумів.}
\answer{ex:08:2}{(а)~Власні числа $-(1\pm g\beta)/\tau$; різниця загасає за $\tau/(1-g\beta)$: $40\ms$ при $g\beta=0.5$ (різниця входів підсилена втричі), $200\ms$ при $0.9$ (у $19$ разів). (б)~$0.905$ і $1.105$; між межами перемагає лише сильний вхід.}
\answer{ex:08:3}{$P(k)=e^{gu_k/\sigma}/\sum_le^{gu_l/\sigma}$, тобто~\eqref{eq:softmax}; $g=10\ln9\approx22$.}
\answer{ex:08:4}{$\bar\Delta=1/3$, $g^*\approx3\ln(1/h)$: $21$, $14$, $9$ проти $16$--$23$, $11$ і $8$ у моделі.}
\answer{ex:08:5}{$g^*$ від $16$--$23$ при $h=0.001$ до $6$--$8$ при $h=0.2$; оцінка $3\ln(1/h)$ правильна в межах півтора раза при $h\le0.05$; при $h\gtrsim\eta/10$ правило $Q\leftarrow Q+\eta(R-Q)$ не встигає засвоїти здобуте пошуком, і $g^*$ застрягає близько $8$.}
\answer{ex:08:6}{$T_p=\tau\ln9=44\ms$ при $g_{\text{lo}}=0$ (модель $44\ms$) і $70\ms$ при $g_{\text{lo}}=0.25$: пачка в $2$--$3\tau$ за підсилення близько нуля.}
\answer{ex:08:7}{(а)~Найкраще стале $0.925$ ($g=8$), оракул $0.929$ ($16/8$) --- вигравати майже нічого. (б)~Наприклад, спокійні епохи з рідкісними винагородами ($p\in[0,0.3]$) і мінливі з $h=0.1$: $0.850$ проти $0.872$; підлаштування забирає близько чверті при $\eta=0.2$ і майже все при $\eta=0.05$.}
